% ============================================================
% estilo_beamer.tex — preámbulo de la presentación
% Réplica de la paleta de 04-documentos/estilo.tex para que el
% informe y el deck se vean como un mismo trabajo.
% ============================================================

\usepackage[utf8]{inputenc}
\usepackage[T1]{fontenc}
\usepackage[spanish,es-noshorthands]{babel}
\usepackage{lmodern}
\usepackage{booktabs}
\usepackage{array}
\usepackage{colortbl}
\usepackage{tabularx}
\usepackage{graphicx}
\usepackage{xcolor}

% ---------- paleta compartida con el informe ----------
\definecolor{ink}{RGB}{0,62,92}
\definecolor{inklight}{RGB}{92,132,153}
\definecolor{rulegray}{RGB}{205,214,219}
\definecolor{critred}{RGB}{158,38,32}
\definecolor{amber}{RGB}{176,98,16}
\definecolor{slate}{RGB}{104,112,118}
\definecolor{headrow}{RGB}{0,62,92}
\definecolor{softrow}{RGB}{242,246,248}

% ---------- tema ----------
\usetheme{default}
\usecolortheme{default}
\setbeamertemplate{navigation symbols}{}
\setbeamercolor{structure}{fg=ink}
\setbeamercolor{frametitle}{fg=ink}
\setbeamercolor{title}{fg=ink}
\setbeamercolor{itemize item}{fg=inklight}
\setbeamercolor{itemize subitem}{fg=inklight}
\setbeamerfont{frametitle}{series=\bfseries,size=\large}
\setbeamerfont{title}{series=\bfseries}
\setbeamertemplate{itemize item}{\textbf{\textperiodcentered}}

% línea bajo el título de cada lámina
\setbeamertemplate{frametitle}{%
  \vspace{6pt}%
  \usebeamerfont{frametitle}\usebeamercolor[fg]{frametitle}\insertframetitle\par
  \vspace{2pt}{\color{rulegray}\rule{\textwidth}{0.6pt}}\par\vspace{-2pt}%
}

% ---------- numeración de diapositivas (exigida por los lineamientos) ----------
\setbeamertemplate{footline}{%
  \hfill\usebeamerfont{footline}\color{slate}\scriptsize
  \insertframenumber{} / \inserttotalframenumber\hspace{12pt}\vspace{6pt}%
}

% ---------- tablas, igual que en el informe ----------
\newcolumntype{L}[1]{>{\raggedright\arraybackslash}p{#1}}
\newcommand{\thead}[1]{\textcolor{white}{\bfseries #1}}
\renewcommand{\arraystretch}{1.2}
\setlength{\tabcolsep}{4pt}

% ---------- cifra destacada ----------
% #1 = número   #2 = qué es
\newcommand{\cifra}[2]{%
  \begin{minipage}[t]{0.31\textwidth}
    \raggedright
    {\color{ink}\bfseries\fontsize{22}{24}\selectfont #1}\\[2pt]
    {\color{slate}\footnotesize #2}
  \end{minipage}%
}

% variante apilada: ocupa el ancho de su columna
\newcommand{\cifravert}[2]{%
  \par\noindent
  {\color{ink}\bfseries\fontsize{17}{19}\selectfont #1}\ \ {\color{slate}\footnotesize #2}%
  \par\addvspace{6pt}%
}

% ---------- conclusión de la lámina ----------
\newcommand{\remate}[1]{%
  \vspace{4pt}{\color{rulegray}\rule{\textwidth}{0.6pt}}\par\vspace{3pt}
  {\color{ink}\small\textbf{#1}}%
}

% ---------- etiquetas de estado ----------
\newcommand{\etiqueta}[2]{%
  {\setlength{\fboxsep}{2.5pt}\colorbox{#1}{\textcolor{white}{\bfseries\scriptsize\strut\,#2\,}}}%
}
